% GENERATED FILE. Edit canonical lessons, then run npm run generate:books.
% canonical-source-hash: fnv1a64:6db5a5b3448707f2
% canonical-lessons: TE-C532-pempakam, TE-C532-garabam, TE-C532-posincu, TE-C532-bujjagincu, TE-C532-maram

\chapter{Bringing Up Children}
\label{ch:te-bringing-up-children}
\hlchaptermodality{532}

\begin{chapteropening}
\textbf{By the end of this chapter:} \emph{I can talk about bringing up children, coaxing them and their tantrums.}

Complete the last lesson of chapter 532: I can talk about bringing up children, coaxing them and their tantrums.
\end{chapteropening}

\section[\texorpdfstring{\te{పెంపకం} (pempakaṁ)}{పెంపకం (pempakaṁ)}]{\te{పెంపకం} (pempakaṁ) — upbringing}

\label{lesson:TE-C532-pempakam}



\begin{quote}
Before the new one: say the Telugu for an heir, then the Telugu for a lullaby.
\end{quote}

\subsection*{You'll want to know: \te{పెంపకం}}
\textbf{\te{పెంపకం}} — \emph{pempakaṁ} — \textquotedblleft{}upbringing\textquotedblright{}.

\begin{quote}
\te{పిల్లల} \te{మంచి} \te{నడవడికి} \te{తల్లిదండ్రుల} \te{పెంపకమే} \te{కారణం}.
\end{quote}

\begin{quote}
\emph{pillala mañci naḍavaḍiki tallidaṇḍrula pempakamē kāraṇaṁ.}
\end{quote}

\begin{quote}
Good parenting is the reason for the children's good behaviour.
\end{quote}

The noun of \te{పెంచు}, to raise.

\begin{cousinweb}
\textbf{\te{పెంపకం}} (\emph{pempakaṁ}) is native Dravidian: DEDR 4436 lists Telugu \emph{pempakamu}, \textquotedblleft{}rearing\textquotedblright{}, beside \emph{pen(u)cu}, \textquotedblleft{}to nourish, to foster, to rear\textquotedblright{}.
\end{cousinweb}

\begin{grammarlens}[title={What you've built}]
One more word for describing what you lived and felt.
\end{grammarlens}

\subsection*{Your turn}
\begin{itemize}
\raggedright
  \item \emph{Say it:} \emph{pempakaṁ}
  \item \emph{Say it:} \emph{pempakaṁ}, once more
  \item \emph{Say it:} \emph{pillala mañci naḍavaḍiki tallidaṇḍrula pempakamē kāraṇaṁ.} — Good parenting is the reason for the children's good behaviour.
  \item \emph{Recall:} say the Telugu for an heir, then the Telugu for a lullaby, then say \emph{pempakaṁ} again
\end{itemize}

\begin{tcolorbox}[breakable,colback=teal!4,colframe=teal!35!black,title={Before you move on}]
What does \te{పెంపకం} mean? (\textquotedblleft{}Upbringing\textquotedblright{}.) Say it once more.
\end{tcolorbox}
\section[\texorpdfstring{\te{గారాబం} (gārābaṁ)}{గారాబం (gārābaṁ)}]{\te{గారాబం} (gārābaṁ) — pampering; indulgence}

\label{lesson:TE-C532-garabam}



\begin{quote}
Before the new one: say the Telugu for a lullaby, then the Telugu for upbringing.
\end{quote}

\subsection*{You'll want to know: \te{గారాబం}}
\textbf{\te{గారాబం}} — \emph{gārābaṁ} — \textquotedblleft{}pampering; indulgence\textquotedblright{}.

\begin{quote}
\te{అతి} \te{గారాబం} \te{వల్ల} \te{పిల్లలు} \te{మొండిగా} \te{తయారవుతారు}.
\end{quote}

\begin{quote}
\emph{ati gārābaṁ valla pillalu moṇḍigā tayāravutāru.}
\end{quote}

\begin{quote}
Too much pampering makes children stubborn.
\end{quote}

\te{గారాబంగా} \te{పెంచు}, to bring up a child spoilt.

\begin{cousinweb}
\textbf{\te{గారాబం}} (\emph{gārābaṁ}) is a Telugu word for doting indulgence; its origin is not settled.
\end{cousinweb}

\begin{grammarlens}[title={What you've built}]
One more word for describing what you lived and felt.
\end{grammarlens}

\subsection*{Your turn}
\begin{itemize}
\raggedright
  \item \emph{Say it:} \emph{gārābaṁ}
  \item \emph{Say it:} \emph{gārābaṁ}, once more
  \item \emph{Say it:} \emph{ati gārābaṁ valla pillalu moṇḍigā tayāravutāru.} — Too much pampering makes children stubborn.
  \item \emph{Recall:} say the Telugu for a lullaby, then the Telugu for upbringing, then say \emph{gārābaṁ} again
\end{itemize}

\begin{tcolorbox}[breakable,colback=teal!4,colframe=teal!35!black,title={Before you move on}]
What does \te{గారాబం} mean? (\textquotedblleft{}Pampering; indulgence\textquotedblright{}.) Say it once more.
\end{tcolorbox}
\section[\texorpdfstring{\te{పోషించు} (pōṣiñcu)}{పోషించు (pōṣiñcu)}]{\te{పోషించు} (pōṣiñcu) — to support; to provide for}

\label{lesson:TE-C532-posincu}



\begin{quote}
Before the new one: say the Telugu for upbringing, then the Telugu for pampering.
\end{quote}

\subsection*{You'll want to know: \te{పోషించు}}
\textbf{\te{పోషించు}} — \emph{pōṣiñcu} — \textquotedblleft{}to support; to provide for\textquotedblright{}.

\begin{quote}
\te{అన్నయ్యే} \te{కుటుంబాన్ని} \te{పోషించాడు}.
\end{quote}

\begin{quote}
\emph{annayyē kuṭumbānni pōṣiñcāḍu.}
\end{quote}

\begin{quote}
It was my elder brother who provided for the family.
\end{quote}

Past (I …): \te{పోషించాను}, \emph{pōṣiñcānu}.

Also to play a part: \te{పాత్ర} \te{పోషించు}, to play a role.

\begin{cousinweb}
\textbf{\te{పోషించు}} (\emph{pōṣiñcu}) is built on Sanskrit \emph{poṣaṇa}, \textquotedblleft{}nourishing, fostering, supporting\textquotedblright{} (Monier-Williams), with the Telugu verb ending -\te{ఇంచు}.
\end{cousinweb}

\begin{grammarlens}[title={What you've built}]
One more word for describing what you lived and felt.
\end{grammarlens}

\subsection*{Your turn}
\begin{itemize}
\raggedright
  \item \emph{Say it:} \emph{pōṣiñcu}
  \item \emph{Say it:} \emph{pōṣiñcu}, once more
  \item \emph{Say it:} \emph{annayyē kuṭumbānni pōṣiñcāḍu.} — It was my elder brother who provided for the family.
  \item \emph{Recall:} say the Telugu for upbringing, then the Telugu for pampering, then say \emph{pōṣiñcu} again
\end{itemize}

\begin{tcolorbox}[breakable,colback=teal!4,colframe=teal!35!black,title={Before you move on}]
What does \te{పోషించు} mean? (\textquotedblleft{}To support; to provide for\textquotedblright{}.) Say it once more.
\end{tcolorbox}
\section[\texorpdfstring{\te{బుజ్జగించు} (bujjagiñcu)}{బుజ్జగించు (bujjagiñcu)}]{\te{బుజ్జగించు} (bujjagiñcu) — to coax; to cajole}

\label{lesson:TE-C532-bujjagincu}



\begin{quote}
Before the new one: say the Telugu for pampering, then the Telugu for to support.
\end{quote}

\subsection*{You'll want to know: \te{బుజ్జగించు}}
\textbf{\te{బుజ్జగించు}} — \emph{bujjagiñcu} — \textquotedblleft{}to coax; to cajole\textquotedblright{}.

\begin{quote}
\te{మందు} \te{తాగమని} \te{అమ్మ} \te{పాపను} \te{బుజ్జగించింది}.
\end{quote}

\begin{quote}
\emph{mandu tāgamani amma pāpanu bujjagiñcindi.}
\end{quote}

\begin{quote}
Mum coaxed the little girl to take her medicine.
\end{quote}

Past (I …): \te{బుజ్జగించాను}, \emph{bujjagiñcānu}.

Gentle persuasion, usually of a child.

\begin{cousinweb}
\textbf{\te{బుజ్జగించు}} (\emph{bujjagiñcu}) is a Telugu verb for coaxing a child; its origin is not settled.
\end{cousinweb}

\begin{grammarlens}[title={What you've built}]
One more word for describing what you lived and felt.
\end{grammarlens}

\subsection*{Your turn}
\begin{itemize}
\raggedright
  \item \emph{Say it:} \emph{bujjagiñcu}
  \item \emph{Say it:} \emph{bujjagiñcu}, once more
  \item \emph{Say it:} \emph{mandu tāgamani amma pāpanu bujjagiñcindi.} — Mum coaxed the little girl to take her medicine.
  \item \emph{Recall:} say the Telugu for pampering, then the Telugu for to support, then say \emph{bujjagiñcu} again
\end{itemize}

\begin{tcolorbox}[breakable,colback=teal!4,colframe=teal!35!black,title={Before you move on}]
What does \te{బుజ్జగించు} mean? (\textquotedblleft{}To coax; to cajole\textquotedblright{}.) Say it once more.
\end{tcolorbox}
\section[\texorpdfstring{\te{మారాం} (mārāṁ)}{మారాం (mārāṁ)}]{\te{మారాం} (mārāṁ) — a tantrum; whining}

\label{lesson:TE-C532-maram}



\begin{quote}
Before the new one: say the Telugu for to support, then the Telugu for to coax.
\end{quote}

\subsection*{You'll want to know: \te{మారాం}}
\textbf{\te{మారాం}} — \emph{mārāṁ} — \textquotedblleft{}a tantrum; whining\textquotedblright{}.

\begin{quote}
\te{బొమ్మ} \te{కోసం} \te{పాప} \te{మారాం} \te{చేసింది}.
\end{quote}

\begin{quote}
\emph{bomma kōsaṁ pāpa mārāṁ cēsindi.}
\end{quote}

\begin{quote}
The little girl threw a tantrum for a toy.
\end{quote}

\te{మారాం} \te{చేయు}, to whine and sulk to get something, as small children do.

\begin{cousinweb}
\textbf{\te{మారాం}} (\emph{mārāṁ}) is a Telugu word for a child's whining tantrum; its origin is not settled.
\end{cousinweb}

\begin{grammarlens}[title={What you've built}]
One more word for describing what you lived and felt.
\end{grammarlens}

\subsection*{Your turn}
\begin{itemize}
\raggedright
  \item \emph{Say it:} \emph{mārāṁ}
  \item \emph{Say it:} \emph{mārāṁ}, once more
  \item \emph{Say it:} \emph{bomma kōsaṁ pāpa mārāṁ cēsindi.} — The little girl threw a tantrum for a toy.
  \item \emph{Recall:} say the Telugu for to support, then the Telugu for to coax, then say \emph{mārāṁ} again
\end{itemize}

\begin{tcolorbox}[breakable,colback=teal!4,colframe=teal!35!black,title={Before you move on}]
What does \te{మారాం} mean? (\textquotedblleft{}A tantrum; whining\textquotedblright{}.) Say it once more.
\end{tcolorbox}
